% Fragmento público generado desde propietarios íntegros.
% Residencia: 52.13.
% Fuente: ../incoming/radion_monografia_20260827/manuscrito/sections/13_sintesis.tex:1-80
\noindent La síntesis operatoria reúne los lectores de fase, acción, profundidad y memoria sin identificarlos entre sí.\par\medskip

\subsubsection{Cuádruple realización del estado radional}

Sea \(x_\infty\in X_{\mathrm{enr}}\) la sección coinductiva. El radión completo no
es un punto aislado, sino el paquete
\begin{equation}
\mathfrak R_\sigma(x_\infty)
=\bigl(
\theta=1,\sigma,
\pih,\alphah,A,C^*,\Delta_4,
\Phi_T,D_E,\epsone,\hret,
q_{\mathrm{mem}},\mathcal I_\infty
\bigr).
\label{eq:full-radion}
\end{equation}
Cuatro lectores extraen de él objetos diferentes:

\begin{figure}[H]
\centering
\begin{tikzpicture}[
  node distance=17mm and 18mm,
  state/.style={draw=RadNavy,very thick,rounded corners,fill=RadCream,
    minimum width=39mm,minimum height=13mm,align=center},
  readout/.style={draw=RadTeal,rounded corners,fill=RadMist,
    minimum width=39mm,minimum height=12mm,align=center},
  arr/.style={-{Latex},thick,RadCopper}]
\node[state] (x) {estado enriquecido\\\(x_\infty\)};
\node[readout,above left=of x] (circ) {círculo \(S^1\)\\fase};
\node[readout,above right=of x] (ell) {elipse positiva\\acción y forma};
\node[readout,below left=of x] (klein) {interior Klein\\profundidad};
\node[readout,below right=of x] (helix) {hélice nonádica\\memoria y retorno};
\draw[arr] (x)--(circ) node[midway,below left]{\small \(P_{\mathrm{ph}}\)};
\draw[arr] (x)--(ell) node[midway,below right]{\small \(P_{\mathrm{act}}\)};
\draw[arr] (x)--(klein) node[midway,above left]{\small \(P_K\)};
\draw[arr] (x)--(helix) node[midway,above right]{\small \(P_{\mathrm{mem}}\)};
\end{tikzpicture}
\caption{Cuatro publicaciones tipadas del mismo antecedente. La unidad del
estado no elimina la diferencia de codominios.}
\label{fig:four-readers}
\end{figure}

El círculo conserva la clase de fase y olvida el lift. La elipse conserva
acción, anisotropía y orientación simpléctica. Klein dota al interior de una
distancia proyectiva. La hélice conserva número de retornos y refinamiento.
La Mecánica Dimensional no consiste en mezclarlos, sino en construir los
transductores que los hacen compatibles.

\subsubsection{Interpretación operatoria del cociclo \(\varepsilon_R\)}

\(\varepsilon_R\) no es una constante universal independiente de todo contexto. Es la
coordenada de un cociclo de transición asociado a la carta del radión. Su
dominio contiene dos secciones del mismo estado:
\[
S_E=1+D_E,
\qquad
S_T=1+\Phi_T.
\]
Su regla es
\begin{equation}
\epsone=S_T-S_E=\Phi_T-D_E.
\label{eq:epsilon-final-meaning}
\end{equation}
Su codominio es la recta adimensional de unidad angular; la carta
\(360/T_+\) la publica en grados.

\begin{exactbox}
\textbf{En lenguaje llano:} la sección directa y la sección torsional llegan
a la misma celda de unidad por dos rutas internas diferentes. Sus partes
enteras coinciden; sus restos no. \(\epsR\) es la diferencia orientada de esos
restos, calculada con el acarreo de APP y llevada al límite por los cilindros
TPK. No se introduce para que el radián cierre: aparece antes de consultar el
radián y, cuando se publica, lo cierra.
\end{exactbox}

Esta definición explica por qué el número es pequeño. No porque se haya
buscado una corrección pequeña, sino porque \(S_E\) y \(S_T\) comparten un
prefijo profundo. La pequeñez mide la cercanía estructural de dos
publicaciones; el signo mide su orientación; las cifras registran el acarreo.

